\documentclass[11pt]{article}
\usepackage[T1]{fontenc}
\usepackage{lmodern}
\usepackage[margin=1in]{geometry}
\usepackage{amsmath,amssymb,amsthm,microtype}
\setlength{\emergencystretch}{2em}
\usepackage[colorlinks=true,linkcolor=blue,citecolor=blue,urlcolor=blue]{hyperref}
\hypersetup{pdftitle={A sharp dyadic threshold in a characteristic-two random-cone construction},pdfauthor={Alec Kriebel},pdfsubject={Model-specific spectral calibration and attributed group-ring consequences}}
\newtheorem{theorem}{Theorem}
\newtheorem{proposition}[theorem]{Proposition}
\newcommand{\Ftwo}{\mathbb F_2}
\newcommand{\KG}{K[G]}
\title{A sharp dyadic threshold in a\\characteristic-two random-cone construction}
\author{Alec Kriebel\\\small ORCID: \href{https://orcid.org/0009-0001-9320-500X}{0009-0001-9320-500X}}
\date{October 6, 2026}
\begin{document}
\maketitle
\begin{abstract}
OpenAI's October 4 construction uses a projective plane of order
$q=128$ and seven extra letters to obtain a finitely presented torsion-free
group with a finite two-dimensional classifying complex and a
characteristic-two group algebra that is not directly finite.
We determine the least dyadic order $q\ge4$ at which its unchanged squared
turn weights have exponential word-mass decay: it is $q=32$.
A three-class reduction and exact rational certificates prove contraction
at $32$ and exponential growth at $4,8,16$. Checking the remaining
parameter-dependent estimates yields the same construction on a
532-edge rose. This is a threshold for this specified model, rather than
an optimum among all group-algebra examples. Classical HNN embedding
also gives a two-generator realization. We spell out the inherited
scalar idempotent and nonzero absorbed cyclic projective module, which
remain valid over every characteristic-two field. The construction is
existential; no numerical matching or group-ring multiplication
certificate is claimed.
\end{abstract}

\section{Scope and inherited input}
For a torsion-free group $G$ and field $K$, the scalar idempotent
conjecture asks whether $e^2=e$ in $K[G]$ forces $e\in\{0,1\}$
\cite[Problem 1(c)]{Oinert}. The positive-characteristic version includes
characteristic two. The pivotal input here is the exact October 4
manuscript \cite[Theorem 1.1]{OpenAI}, at repository commit
\href{https://github.com/openai/math/tree/adc7f1241b42e322a6451854ab7e4b4c146bf78a}{\texttt{adc7f124}}:
it constructs a finitely presented torsion-free group with a finite
2-dimensional $K(G,1)$ and scalar elements
\begin{equation}\label{eq:input}
ab=1,\qquad ac=0,\qquad c\ne0\qquad\text{in }\Ftwo[G].
\end{equation}
The breakthrough and construction belong to OpenAI. Its probabilistic
and geometric framework also draws on its September 23 zero-divisor
manuscript \cite{OpenAIZD}. Our independent source audits found no
substantive gap in the arguments used; their precise scope is recorded
in the supplement. The earlier family-197 Lean example has torsion and
does not formalize this October 4 theorem. No full formal verification
or conventional human refereeing is claimed.

Our contribution is the exact dyadic threshold for the \emph{unchanged}
seven-extra model in Section~\ref{sec:threshold}, together with the
parameter check in Section~\ref{sec:adaptation}. The defect-idempotent
and module consequences were already publicly recorded in the author's
triage \cite{Triage}; they are included for interpretation, not offered
as a new discovery. The two-generator specialization in
Section~\ref{sec:embedding} uses classical machinery. We make no
first-priority claim and no characteristic-zero, odd-characteristic or
reduced group $C^*$-algebra projection claim.

\section{The exact squared-word threshold}\label{sec:threshold}
Let $q=2^k$ with $k\ge2$, $v=q^2+q+1$ and $p=(q+1)/v$.
The source's signed alphabet $T$ consists of the $v$ points of
$\operatorname{PG}(2,q)$ and seven extra letters. Pair each extra with a
distinct ordinary letter, and pair the remaining $v-7$ ordinary letters
among themselves; write $\bar t$ for the inverse of $t$. For a reduced
turn $u\ne\bar t$, retain exactly the source weights
\[
w(t,u)=\begin{cases}
(q+1)^{-1},&\bar t,u\text{ both ordinary},\\
p,&\bar t,u\text{ of different kinds},\\
1/2,&\bar t,u\text{ both extra},
\end{cases}\qquad w(t,\bar t)=0.
\]
For $W=t_1\cdots t_h$, set $P(W)=\prod_{i=1}^{h-1}w(t_i,t_{i+1})$.
These are turn-frequency weights, not probabilities assigned to words.

\begin{theorem}\label{thm:threshold}
In this model, for any allowed inverse pairing, $q=32$ is the least
dyadic $q\ge4$ for which $M=(w(t,u)^2)$ has spectral radius below one.
At $q=32$, for every $h\ge1$,
\[
\sum_{W\in T^h}P(W)^2\le\frac{5376}{5}
       \left(\frac{987}{1000}\right)^{h-1}.
\]
For each $q\in\{4,8,16\}$ the same sum is at least
$C_q(503/500)^{h-1}$, where $C_q>0$ is explicit below.
\end{theorem}
\begin{proof}
Partition the actual letters into $E$ (the seven extras), $S$ (their
seven ordinary mates), and $O$ (the remaining $v-7$ ordinary letters).
Vectors constant on these classes are invariant under $M$. In class
order $(E,S,O)$ their exact quotient is
\begin{equation}\label{eq:quotient}
Q(q)=\begin{pmatrix}
7p^2&6/(q+1)^2&(v-7)/(q+1)^2\\
3/2&7p^2&(v-7)p^2\\
7p^2&7/(q+1)^2&(v-8)/(q+1)^2
\end{pmatrix}.
\end{equation}
For a row in $E$, its inverse excludes one $S$ successor; for a row in
$S$, its inverse excludes one $E$ successor; for a row in $O$, its
inverse excludes one $O$ successor. This proves the quotient formula
for every permitted pairing. $Q$ is strictly positive. Its positive
Perron eigenvector lifts to a positive eigenvector of $M$, so
$\rho(M)=\rho(Q)$. The following certificates also prove the bounds
directly, without numerical eigenvalues.

At $q=32$, lift $f=(1,13/5,1)$. The three ratios $(Qf)_i/f_i$ are
\begin{equation}\label{eq:upper}
\frac{95146189}{96562235},\qquad
\frac{579417}{592826},\qquad
\frac{857592557}{869060115},
\end{equation}
each strictly below $\lambda=987/1000$. Hence
$Mf\le\lambda f$. As $\mathbf1\le f$ and
$\mathbf1^{\mathsf T}f=5376/5$, nonnegativity gives
\[
\mathbf1^{\mathsf T}M^{h-1}\mathbf1
 \le\lambda^{h-1}\mathbf1^{\mathsf T}f.
\]
The left side is exactly the stated sum over words.
Thus the source's decay lemma holds, for example with
$\delta=-\frac14\log\lambda>0$ for all sufficiently large $h$.

For $q=4,8,16$, use $x=(193/500,1,97/250)$, with $x\le\mathbf1$.
The exact ratios $(Qx)_i/x_i$ are
\[
\begin{array}{c|ccc}
q&E&S&O\\\hline
4&480733/303975&26959/21000&500011/305550\\
8&10342603/9256473&5726739/5329000&10523789/9304434\\
16&35146373/34932807&3571543/3549000&2078707/2065518
\end{array}
\]
Every entry exceeds $\mu=503/500>1$. Therefore $Mx\ge\mu x$ and
\[
\mathbf1^{\mathsf T}M^{h-1}\mathbf1
 \ge\mu^{h-1}\mathbf1^{\mathsf T}x,
\qquad C_q=7\frac{193}{500}+7+(v-7)\frac{97}{250}>0.
\]
This proves $\rho(M)>1$ for $4,8,16$, completing the least-order claim.
\end{proof}

The lower bound excludes the source's squared-word decay statement at
these orders. It does not prove nonexistence of good matchings or
idempotents there, or exclude a modified type model or another method.
Order $2$ is outside the theorem's domain: the unchanged extra-count
prescription would require both $3m-1$ and $3(m-1)$ divisible by four,
which is impossible. Exact rational verification is supplied as
\texttt{verification/verify\_parameter.py}.

\section{Carrying the construction to order 32}\label{sec:adaptation}
\begin{proposition}\label{prop:smaller}
The October 4 random-cone construction remains valid at $q=32$ with the
same seven-extra types. It gives a finitely presented torsion-free group
$G$ generated by at most $532$ elements, a finite 2-dimensional $K(G,1)$,
and the scalar triple \eqref{eq:input}.
\end{proposition}
\begin{proof}
Here $v=1057$, $p=33/1057$ and $|T|=1064$. For sufficiently large
$m\equiv1\pmod4$, take $m$ vertices in every line class on side $A$
and $m-1$ on side $B$, with $n=vm$. Each ordinary outgoing set is its
33-point line. Give the distinguished $A$ vertex all seven extras;
for each complement of a Fano line assign that four-element extra set
at $(33m-1)/4$ vertices of $A$ and $33(m-1)/4$ of $B$, evenly among
line classes. All other extra sets are empty.
The line-capacity coefficient is $7\cdot33/(4\cdot1057)=33/604<1$,
so disjoint assignments exist for all sufficiently large $m$.
Each extra is in four complements and each extra pair in two, while
ordinary lines have odd intersection size. Thus the source's exact
inverse-slot balances and parity identities persist. Degrees are now
$33,37,40$, with minimum $33$ and maximum $d_*=40$.

The decay input is Theorem~\ref{thm:threshold}. The remaining estimates
in \cite[Sections 3--5]{OpenAI} are verified as follows.
Choose a fixed $c_0>0$ satisfying
$c_0\log(2|T|/p)<1$ and $2c_0\log d_*<1$; $c_0=1/100$ works.
The short-cycle switching proof then has the same error
$r_n=O(\log n+n^{2c_0\log d_*})=o(n)$ and supplies the same conditioned
edge-prescription bounds. It uses only fixed $T,p,d_*$, these strict
inequalities, and exact slot balance, all of which hold here.

For the diameter argument, a $k$-vertex set without factor-two expansion
forces at least $r=\lceil33k/2\rceil$ prescribed edges in a $2k$-vertex
set, with $16.5k\le r\le17k$. The source's union bound becomes
\[
\Pr(\text{failure at }k)
 \le\bigl(C_3(k/n)^{29/2}\bigr)^k.
\]
A fixed small $\eta>0$ can satisfy $17\eta<p/4$ and
$C_3\eta^{29/2}<1/2$. Summing separately over $k\le\sqrt n$ and
$\sqrt n<k\le\eta n$ gives $o(1)$; the first sum uses exponent
$29/4>0$. The ball-packing argument then gives the same $dL$ component
diameter bound, where $L=\lfloor c_0\log n\rfloor$.
The supplemental exact certificate even permits $\eta=10^{-6}$ with
an explicit conservative bound for $C_3$.

All bounded-path estimates in Section 4 of the source use the fixed
alphabet, the positive minimum turn weight, the just-established
prescription bound and exponential squared-word decay. Their constants
may change but their order of quantifiers does not: choose the unpaired
fraction first, then any fixed path-count, length and comparison bounds.
The multiplicity-stage first moments do not require independence between
stages. In the planar reduction, deleting at most two incident edges
leaves degree at least $31\ge3$, so the short-closure lemma is unchanged;
the Euler characteristic and separator estimates contain no further
projective-plane order requirement. Consequently the probability of a
reduced spherical arrangement tends to zero in this conditioned model.
These statements follow by applying the cited proofs with the displayed
replacements, rather than treating their originally fixed-$128$
theorems as unproved general-parameter theorems.

Select one sufficiently large matching outcome with these properties.
The deterministic cone criterion in \cite[Section 6]{OpenAI} uses no
$q$ parameter: it gives asphericity and protects the two roots.
The parity proposition in its Section 2 yields \eqref{eq:input}.
Coning all graph components onto the 532-edge rose yields a finite
2-dimensional classifying complex. It is homotopy equivalent to the
presentation complex on the rose edges with one relator for each
non-tree graph edge, so $G$ is generated by at most 532 elements.
Torsion-freeness follows as in the source from its finite-dimensional
classifying complex.
\end{proof}

For completeness, after choosing an outcome, choose a root and spanning
tree in every component and let $w_x$ be the tree-path word to $x$.
For each oriented non-tree edge $x\xrightarrow{t}y$ impose
$w_xtw_y^{-1}$. At the designated-root components $A',B'$, the witnesses
are $a=\sum_{x\in A'}w_x$, $b=\sum_{x\in A'}w_x^{-1}$ and
$c=\sum_{y\in B'}w_y^{-1}$. Product-graph parity gives $ab=1,ac=0$;
root protection makes the identity coefficient of $c$ equal to one.
The order-32 incidence plane and inverse pairing have a reproducible
finite certificate in the supplement. No numerical $m$ satisfying all
asymptotic requirements, matching outcome, relator list, reduced support
of $e$, or target group-ring multiplication certificate is supplied.

\section{Classical two-generator realization}\label{sec:embedding}
The number 532 describes this random-cone recipe, not the smallest
possible group generating set. We record explicitly why classical HNN
embedding \cite{HNN,Morozov} already yields two generators while
preserving the stronger finite classifying-space condition.

\begin{proposition}\label{prop:embedding}
Every group $G$ with a finite 2-dimensional $K(G,1)$ embeds in a group
$H$ generated by two elements with a finite 2-dimensional $K(H,1)$.
If $G$ has an aspherical presentation on $g_1,\ldots,g_s$ with $r$
relators, $H$ has an aspherical presentation on two generators with
$r$ relators. Torsion-freeness is preserved.
\end{proposition}
\begin{proof}
Set $B=G*F(x,y)$, and in $B$ define
\[
u_0=y,\quad v_0=x,\qquad
u_i=x^{-i}yx^i,\quad v_i=y^{-i}xg_i y^i\quad(1\le i\le s).
\]
The $u_i$ freely generate: they are a finite subset of the standard
free basis of the kernel of $F(x,y)\to\mathbb Z$, $x\mapsto1,y\mapsto0$.
Projection killing $G$ sends the $v_i$ to the corresponding free-basis
subset with $x,y$ interchanged, and proves they freely generate too.
The HNN extension with relations $t^{-1}u_it=v_i$ therefore embeds $B$.
Its zeroth relation gives $y=txt^{-1}$; the others give
$g_i=x^{-1}y^it^{-1}x^{-i}yx^ity^{-i}$. Substituting for $y$ expresses
$g_i$ by the word
\[
W_i=x^{-1}t x^i t^{-2}x^{-i}t x t^{-1}x^i t^2 x^{-i}t^{-1}.
\]
Thus $H=\langle x,t\rangle$, and its presentation is
$\langle x,t\mid R_j(W_1,\ldots,W_s),\ 1\le j\le r\rangle$.

Let $X$ be the finite classifying complex for $G$, with its maximal
1-skeleton tree collapsed. The vertex space $X\vee R_2$ is a finite
$K(B,1)$. Attach the mapping cylinder of a rose with $s+1$ petals via
the two injective maps representing $u_i$ and $v_i$. This gives a finite
complex of dimension at most two for $H$. Its universal cover is a
tree of contractible vertex and edge spaces; successively adjoining
finite subtrees gives contractible homotopy pushouts. Sphere maps have
compact image contained in a finite subcomplex, and hence a finite
subtree, so the cover is contractible. This proves asphericity.
Finite-order elements of an HNN extension are conjugate into its base;
the base free product is torsion-free when $G$ is, proving the assertion.

To see that the final two-generator presentation itself is aspherical,
first consider the subcomplex having only the $s+1$ HNN defining cells.
Each $g_i$ occurs once in its own such cell and in no other such cell,
so those generator-cell pairs can be collapsed there; then collapse
$y$ with the zeroth defining cell. The resulting retraction to the
$(x,t)$ rose sends $g_i$ to $W_i$. Homotopy invariance of attaching
cells carries the original $r$ relators to $R_j(W)$, proving homotopy
equivalence with the displayed presentation complex. This uses these
specific eliminations, not arbitrary Tietze moves.
\end{proof}

An embedding $G\hookrightarrow H$ induces an injective map $K[G]\to K[H]$
because distinct group elements remain distinct basis elements. The
scalar triple is preserved. No support is shortened or computed by this
transport. A torsion-free group generated by at most one element is
trivial or infinite cyclic; its group algebra over a field is a domain.
Consequently two is the least possible number of generators for such
an idempotent example, but the embedding mechanism is classical.

\section{Ring and right-module proof}
The reduction to a scalar idempotent is classical
\cite[Proposition \texttt{kaplansky\_relations}]{Gardam}; the splitting
is the regular right-module interpretation of direct finiteness
\cite[p.~9]{Ara}. We include the maps to fix handedness.
\begin{proposition}\label{prop:ring}
Let $R$ be a nonzero unital ring and suppose \eqref{eq:input} holds in $R$.
Then $e=1-ba$ is an idempotent other than $0$ and $1$, and $P=eR$ has all
the right-module and $K_0$ properties in the consequence theorem below.
No characteristic assumption is needed for this proposition.
\end{proposition}
\begin{proof}
Put $f=ba$. Then $f^2=b(ab)a=f$, so
$e^2=(1-f)^2=1-2f+f^2=1-f=e$.
Also $ec=c-b(ac)=c\ne0$, whence $e\ne0$.
If $e=1$, then $ba=0$, and $a=a(ba)=0$ contradicts $ab=1$ in a
nonzero ring. Thus $e\ne1$.

The complementary idempotents $f$ and $e$ satisfy $fe=ef=0$ and
$f+e=1$. Therefore $R=fR+eR$. If $x=fr=es$, then
$x=fx=fes=0$, giving the direct sum
\begin{equation}\label{eq:split}
R_R=baR\oplus eR.
\end{equation}
Since $baR\subseteq bR$ and $b=bab$, the reverse inclusion holds as
well: $baR=bR$. Left multiplication by $b$ is a right $R$-linear
isomorphism $R_R\to bR$, with inverse left multiplication by $a$.
Indeed $abr=r$ for every $r$, and $ba(br)=br$ for every $br\in bR$.

Consequently \eqref{eq:split} identifies $R_R$ with $R_R\oplus P$.
For a fully explicit pair of inverse maps, use
\[
\Phi(r)=(ar,er),\qquad \Psi(u,p)=bu+p
\quad (r,u\in R,\ p\in eR).
\]
The identities $ae=eb=0$, $ep=p$, and $ab=1$ give
$\Psi\Phi(r)=bar+er=r$ and $\Phi\Psi(u,p)=(u,p)$.
They also show $\ker(r\mapsto ar)=eR$; in particular $c\in P$.

The module $P$ is a direct summand of the free right module $R_R$, hence
projective. It is cyclic, generated by $e$, and nonzero because it
contains $e=e\cdot1\ne0$. By definition, $K_0(R)$ is the group completion
of the direct-sum monoid of finitely generated projective right modules
\cite[Chapter II, Section 2]{Weibel}. Hence the isomorphism implies
$[R]=[R]+[P]$ and therefore $[P]=0$ in that abelian group.
It does not imply $P=0$. Rather, $R\oplus0\cong R\oplus P$ with
$0\not\cong P$ exhibits failure of projective cancellation.
\end{proof}

\begin{theorem}\label{thm:consequence}
There is a two-generator finitely presented torsion-free group $H$ with a
finite 2-dimensional $K(H,1)$ and a scalar $e\in\Ftwo[H]$, $e^2=e$,
$e\ne0,1$. For every characteristic-two field $K$, its same image is
nontrivial in $S=K[H]$ and $P=eS$ is a nonzero cyclic projective right
module satisfying $S_S\cong S_S\oplus P$ and $[P]=0$ in $K_0(S)$.
\end{theorem}
\begin{proof}
Apply Propositions~\ref{prop:smaller} and \ref{prop:embedding} and
transport $a,b,c$ to $\Ftwo[H]$. Proposition~\ref{prop:ring} gives
$e=1-ba$ and the right-module properties.
The unique unital embedding $\Ftwo\hookrightarrow K$ induces an injective
ring map $\Ftwo[H]\to K[H]$: the elements of $H$ are a basis over each
field. Thus $e$ and $e-1$ remain nonzero, and all scalar identities
persist. Apply the ring proposition again over $K[H]$.
\end{proof}

Vanishing of this one class does not assert $K_0(R)=0$. In fact the
scalar inclusion $\Ftwo\to R$ and augmentation $R\to\Ftwo$ split off
$K_0(\Ftwo)=\mathbb Z$ from $K_0(R)$. The characteristic-two field
embedding supplies no coefficient map to a field of odd characteristic
or characteristic zero.

\section{Priority, validation and limitations}
The published triage \cite{Triage}, commit
\texttt{f27318d83bd7000ef817957a9a4b3087de28d198}, already states
$e=1-ba$, $P=eR$, $R\cong R\oplus P$ and $[P]=0$, conditional on
checking the October 4 input. Those core consequences are duplicated.
The present note's additional result is Theorem~\ref{thm:threshold}
and the source adaptation in Proposition~\ref{prop:smaller}.
The inspected pinned manuscript fixes $128$; the relevant companions
and triage do not supply this exact calibration or its lower-order
obstruction. A finite literature search cannot establish universal
priority, and we assert neither a first solution of a named conjecture
nor a new embedding or idempotent mechanism.

The supplement retains separate audits of the actual probabilistic and
cone-picture proofs, a parameter-specific adversarial audit, primary
source identities, exact rational checks, the finite incidence model,
and the embedding word calculation. They certify their stated portions;
finite computations do not certify an unlisted matching or replace the
proof of group existence. No characteristic-zero or odd-characteristic
conclusion, matrix-only substitute, nonzero $K_0$ class or full Lean
formalization is inferred.

\paragraph{Disclosure.}
AI tools were used extensively in research, drafting and verification,
including independent automated adversarial audits. This preprint has
not undergone conventional human peer review or refereeing. OpenAI is
the supplied author of the upstream AI-produced construction; no
affiliation or coauthor is inferred.

\begin{thebibliography}{99}
\raggedright
\bibitem{OpenAI} OpenAI, \emph{A Torsion-Free Group Algebra That Is Not
Directly Finite}, OpenAI Math Release preprint, October 4, 2026,
Theorem 1.1. Pinned source commit
\href{https://github.com/openai/math/tree/adc7f1241b42e322a6451854ab7e4b4c146bf78a/preprints/A-Torsion-Free-Group-Algebra-That-Is-Not-Directly-Finite-October-4-2026}{adc7f124}.
\bibitem{Oinert} Johan {\"O}inert, \emph{Units, zero-divisors and idempotents
in rings graded by torsion-free groups},
\href{https://arxiv.org/abs/1904.04847v3}{arXiv:1904.04847v3}, July 20, 2023
(first submitted April 9, 2019).
\bibitem{Gardam} Giles Gardam, \emph{The Kaplansky conjectures}, lecture
notes, source version February 20, 2024,
\href{https://github.com/gilesgardam/lectures/blob/ec842e6a09851e98a38e2a4bea3b273b4a9debb9/kaplansky.tex}{commit ec842e6}, Proposition
\texttt{kaplansky\_relations}.
\bibitem{Ara} P. Ara, K. R. Goodearl, K. C. O'Meara and E. Pardo,
\emph{Separative cancellation for projective modules over exchange rings},
May 1996 report no. 142, p.~9,
\href{https://ir.canterbury.ac.nz/bitstreams/9810640d-0c02-42b1-a9ac-49f0f8bd703b/download}{institutional full text}.
\bibitem{Weibel} Charles A. Weibel, \emph{The K-book: An Introduction to
Algebraic K-theory}, AMS, 2013, Chapter II,
\href{https://sites.math.rutgers.edu/~weibel/Kbook/Kbook.ii.pdf}{author's chapter}.
\bibitem{Triage} Alec Kriebel, \emph{Ten further research paths opened by
the OpenAI mathematics release}, public research triage, October 6, 2026,
\href{https://github.com/AlecKriebel/Math/blob/f27318d83bd7000ef817957a9a4b3087de28d198/openai_followon_batch2_20261006/REPORT.md}{report, Section 4}; accompanying
\href{https://github.com/AlecKriebel/Math/blob/f27318d83bd7000ef817957a9a4b3087de28d198/openai_followon_batch2_20261006/agent_notes/algebra_groups.md}{algebra notes, Section 1}.
\bibitem{OpenAIZD} OpenAI, \emph{A Torsion-Free Group Algebra with Zero
Divisors}, OpenAI Math Release preprint, September 23, 2026,
\href{https://github.com/openai/math/tree/adc7f1241b42e322a6451854ab7e4b4c146bf78a/preprints/A-Torsion-Free-Group-Algebra-with-Zero-Divisors-September-23-2026}{pinned source}.
\bibitem{HNN} G. Higman, B. H. Neumann and H. Neumann,
\emph{Embedding theorems for groups}, J. London Math. Soc. 24 (1949),
247--254, Theorems I and IV,
\href{https://doi.org/10.1112/jlms/s1-24.4.247}{doi:10.1112/jlms/s1-24.4.247}.
\bibitem{Morozov} A. Morozov and P. Schupp, \emph{Computable permutations
and word problems}, Enseign. Math. 64 (2018), 143--160,
Observation 3.4, p.~151,
\href{https://doi.org/10.4171/LEM/64-1/2-6}{doi:10.4171/LEM/64-1/2-6}.
\end{thebibliography}
\end{document}
